\documentclass[a4paper,11pt]{book}
\usepackage[T1]{fontenc}
\usepackage{lmodern}
\usepackage{babel}
\babelprovide[import,main]{french}

\title{Un petit livre}
\author{L'équipe des tests}
\date{1\textsuperscript{er} janvier 2026}

\begin{document}

\frontmatter
\maketitle
\tableofcontents

\chapter{Avant-propos}

Ce livre essaie la classe \emph{book} avec le paquet babel en français :
les guillemets, les espaces avant la ponctuation haute et les noms des
parties du livre.

\mainmatter

\chapter{Le début}\label{chap:debut}

\section{Les caractères}

Le français a des accents : é, è, ê, à, ù, ç, ï et œ. Il met des
guillemets autour d'une citation : « Bonjour ! » dit-il. Une question
finit par un point d'interrogation ? Une note de bas de page donne un
détail\footnote{Une note courte, en bas de la page.}.

\section{Une liste}

Les étapes du travail :
\begin{enumerate}
  \item lire le fichier ;
  \item composer les pages ;
  \item dessiner les pages.
\end{enumerate}

\chapter{La suite}

Le chapitre~\ref{chap:debut} commence à la page~\pageref{chap:debut}.
Les livres de Knuth~\cite{knuth} et de Lamport~\cite{lamport} décrivent
\TeX{} et \LaTeX.

\begin{thebibliography}{9}
\bibitem{knuth} Donald E. Knuth, \emph{The \TeX book}, Addison-Wesley, 1984.
\bibitem{lamport} Leslie Lamport, \emph{\LaTeX: A Document Preparation
  System}, Addison-Wesley, 1994.
\end{thebibliography}

\end{document}
